\section{Extended Relational Algebra}
\label{sec:prelim}

We now present the query language underlying ProvSQL; it is an extension
of the relational algebra with multiset (or \emph{bag}) semantics, allowing for
aggregation. We first give basic definitions and notation about multisets.

\begin{definition}
	A \emph{multiset} m over a set~$V$ is a function $V\to\mathbb{N}^*$. It
	is
	is \emph{finite} if it has finite domain. For $v\in V$, we note $v\in m$ if $m(v)>0$
	and $v\not\in m$ otherwise.
	Given two
	multisets $m_1$ and $m_2$ over~$V$, we define the \emph{multiset union} of
	$m_1$ and~$m_2$ as $m_1\uplus m_2: x\mapsto
		m_1(x)+m_2(x)$ and the \emph{cross product} of $m_1$ and~$m_2$ as the multiset
	over $V^2$ defined by $m_1\times m_2:(x_1,x_2)\mapsto m_1(x_1)\times
		m_2(x_2)$. Any set~$V$ can be seen as the multiset over~$V$ where $m(v)=1$ for
	all $v\in V$.

	A finite multiset with
	domain $\{a_1,\dots,a_n\}$ can be
	written in the form
	\(\mset{\underbrace{a_1,\dots,a_1}_{\text{$m(a_1)$ times}},\dots,
	\underbrace{a_n,\dots,a_n}_{\text{$m(a_n)$ times}}}.\)
	Finally, we also use the notation $\mset{f(x)\mid x\in m}$ or
	$\biguplus_{x\in m} \mset{f(x)}$ for
	a finite multiset~$m$ with domain $\{a_1,\dots,a_n\}$ and some function
	$f$ over~$V$
	to mean the multiset
	\(\mset{\underbrace{f(a_1),\dots,f(a_1)}_{\text{$m(a_1)$ times}},\dots,
	\underbrace{f(a_n),\dots,f(a_n)}_{\text{$m(a_n)$ times}}}.\)
\end{definition}

We define an algebra for queries over relational databases under the
multiset semantics
which is as close as possible to a simple subset of SQL as
dealt with by standard DBMSs.
Our algebra is based on the relational algebra with set
semantics typically used in database
theory~\cite{DBLP:books/aw/AbiteboulHV95} but includes
operators to explicitly control
duplicate elimination as
in~\cite{DBLP:conf/pods/DayalGK82,DBLP:journals/iandc/GrumbachM99} and
to express aggregation \cite{DBLP:journals/tcs/Libkin03}.

To simplify the presentation, we adopt an ordered, unnamed, and untyped
perspective: attributes are referred to by their positions within the
relation, do not have a name, and are assumed to be from a universal
domain~$\mathcal{V}$. In practical SQL implementations, attributes do have names and
types, but it is straightforward to propagate them throughout query
execution.

% lax begin Lax392996.Databases
Let us fix a set $\mathcal{L}$ of relation labels and a set $\mathcal{V}$
of values. A~database schema
$\mathcal{D}:\mathcal{L}\to\bN$ assigns arities to a \emph{finite}
set of relation labels.
A \emph{relation of arity~$k\in\bN$} is a \emph{finite multiset} of
tuples over~$\mathcal{V}^k$. An instance of a database
schema~$\mathcal{D}$, or a database over~$\mathcal{D}$ for short, is a
function mapping each relation name $R$ in the domain of~$\mathcal{D}$ to
a relation of arity $\mathcal{D}(R)$.
% lax end

A \emph{positional index} is an expression of the form ``$\#i$''
for $i\in\bN^*$. A \emph{term} is any expression involving
positional indices and values from~$\mathcal{V}$, combined using
arbitrary operators and functions over values in $\mathcal{V}$. The
\emph{max-index} of a term is the maximum of all $i$ such that
``\#$i$'' appears in the term (or 0 if none appear). Given a
tuple~$u=(u_1,\dots,u_k)$ and a term~$t$ of max-index~$\leq k$, we
define $t(u)$ as the value of~$\mathcal{V}$ obtained by replacing in~$t$
every
positional index~\#$i$ with~$u_i$ and evaluating the whole
expression.

% lax begin Lax392996.RelationalAlgebra
Given a database schema~$\mathcal{D}$, we recursively define the language
$\RA_k$ of relational algebra queries of arity $k\in\mathbb{N}$ as
follows \lean{Query}{Query}:
\begin{compactdesc}
	\item[relation] for any relation~$R$ in the domain of~$\mathcal{D}$,
	$R\in\RA_{\mathcal{D}(R)}$;
	\item[projection] for $k\in\bN$, $q\in\RA_{k}$,
	$t_1,\dots, t_n$ terms of max-index~\mbox{$\leq k$}, $\Pi_{t_1,\dots,t_n}(q)\in\RA_n$;
	\item[selection] for $k\in\bN$, $q\in\RA_k$, and $\phi$ a
	Boolean combination of (in)equality comparisons between terms
	of max-index~$\leq k$,
	$\sigma_\phi(q)\in\RA_k$;
	\item[cross product] for $k_1,k_2\in\bN$, $q_1\in\RA_{k_1}$,
	$q_2\in\RA_{k_2}$, $q_1\times q_2\in\RA_{k_1+k_2}$;
	\item[multiset sum] for $k\in\bN$, $q_1,q_2\in\RA_k$,
	$q_1\uplus q_2\in\RA_k$;
	\item[duplicate elimination] for $k\in\bN$,
	$q\in\RA_k$, $\epsilon(q)\in\RA_k$;
	\item[multiset difference] for $k\in\bN$, $q_1,q_2\in\RA_k$, $q_1 - q_2\in\RA_k$;
	\item[aggregation] for $k\in\bN$,
	$q\in\RA_k$, distinct $(i_j)_{1\leq j\leq k}$,
	terms~$t_1,\dots t_n$
	of max-index~$\leq k$, functions $f_1,\dots,f_n$ from finite multisets of values to
	values,
	$\gamma_{i_1,\dots,i_m}[t_1:f_1,\dots,t_n:f_n](q)\in\RA_{m+n}$.
\end{compactdesc}
% lax end

Some other operators can be seen as syntactic sugar:
\begin{inparadesc}
	\item[join] $q_1\bowtie_\phi q_2 \defeq \sigma_\phi(q_1\times q_2)$;
	%	\item[full join]
	%	$q_1\bowtie q_2\defeq q_1\bowtie_{\#1=\#(k+1)\land\dots\land\#k=\#(2k)}q_2\:$
	%	(if $q_1,q_2\in\RA_k$);
	\item[set union] $q_1\cup q_2\defeq \epsilon(q_1\uplus q_2)$.
\end{inparadesc}

\begin{toappendix}
% lax begin Lax392996.MultisetSemantics
We define the semantics $\sem{I}{\cdot}$ of relational algebra
queries on an instance~$I$ over~$\mathcal{D}$ by induction:
\lean{Query}{Query.evaluate}
\begin{compactdesc}
	\item[relation] for any relation label~$R$ in the domain of~$\mathcal{D}$,
	$\sem{I}{ R}\defeq I(R)$;
	\item[projection] for $k\in\bN$, $q\in\RA_{k}$,
	$t_1,\dots, t_n$ terms of max-index~\mbox{$\leq k$},
	$\sem{I}{\Pi_{t_1,\dots,t_n}(q)}\defeq
		\mset{(t_1(u),\dots,t_n(u))\mid u\in\sem{I}{ q}}
	$;
	\item[selection] for $k\in\bN$, $q\in\RA_k$, and $\phi$ a
	Boolean combination of (in)equality comparisons between terms
	of max-index~$\leq k$,
	$\sem{I}{\sigma_\phi(q)}\defeq
		\mset{u\mid u\in\sem{I}{ q}, \phi(u)}$;
	\item[cross product] for $k_1,k_2\in\bN$, $q_1\in\RA_{k_1}$,
	$q_2\in\RA_{k_2}$, \(\sem{I}{ q_1\times q_2} \defeq
	\sem{I}{q_1}\times\sem{I}{q_2};\)
	\item[multiset sum] for $k\in\bN$, $q_1,q_2\in\RA_k$,
	\(\sem{I}{ q_1\uplus q_2}\defeq\sem{I}{ q_1}\uplus\sem{I}{ q_2};\)
	\item[duplicate elimination] for $k\in\bN$,
	$q\in\RA_k$,
	\(
	\sem{I}{\epsilon(q)}\) maps $t$ to~$1$
	if $\sem Iq(t)>0$ and to~$0$ otherwise;
	\item[multiset difference]
	for $k\in\bN$, $q_1,q_2\in\RA_k$,
	$\sem{I}{q_1-q_2}$
	maps $t$ to~$0$ if
	$\sem{I}{q_2}(t)>0$ and to~$\sem{I}{q_1}(t)$ otherwise;
	\item[aggregation] for $k\in\bN$,
	$q\in\RA_k$, distinct $(i_j)_{1\leq j\leq k}$,
	terms~$t_1,\dots t_n$
	of max-index~$\leq k$, functions $f_1,\dots,f_n$ from finite multiset of values to
	values,
	\(\sem{I}{\gamma_{i_1,\dots,i_m}
		[t_1:f_1,\dots,t_n:f_n](q)}\) is:\\
    \(
	\Big\{\,
	\Big(v_1,\dots,v_m,                f_1\left(\mset{t_1(u)\mid u\in \sem{I}{q},
	(u_{i_1},\dots,u_{i_m})=(v_1,\dots,v_m)}\right),
  \dots,\)\\{}\qquad
	\(f_n\left(\mset{t_n(u)\mid
	u\in
	\sem{I}{q},
	(u_{i_1},\dots,u_{i_m})=(v_1,\dots,v_m)}\right)\Big)
	\:\Big|
	(v_1,\dots,v_m)\in\sem{I}{\epsilon(\Pi_{\#i_1,\dots,\#i_m}(q))}
	\,\Big\}.
	\)
  \bigskip
\end{compactdesc}
% lax end
\end{toappendix}

\revision{The semantics $\sem{I}{\cdot}$ of queries on an instance
  $I$ over~$\mathcal{D}$ is defined as expected (see companion repository
for details) for most operators.}
Note\revision{, however,} that the definition of the multiset difference is not the usual one,
and also not the one corresponding to the \mintinline{sql}{EXCEPT ALL} operator of
SQL, though it does correspond to the important \mintinline{sql}{NOT IN}
operator of SQL. This is mainly for practicality: the standard definition of multiset
difference
where $\sem{I}{q_1-q_2}(t)=\max(0,\sem{I}{q_1}(t)-\sem{I}{q_2}(t))$ leads
to intractability of the provenance computation. The difference
disappears in the set semantics: the semantics of $\epsilon(q_1-q_2)$
matches with that of \mintinline{sql}{EXCEPT}.

Importantly, in this work, we assume that the aggregation operator, if it is used, is
applied last (at the top-level operator), similarly to what is done in
Section~3 of~\cite{amsterdamer2011provenance}; nested aggregation
queries, discussed in Section~4 of~\cite{amsterdamer2011provenance}, are
implemented in ProvSQL, but as there is currently no clear semantics for
semiring provenance for arbitrary semirings over nested aggregation
queries, we leave them out-of-scope of this work.

% lax begin Lax392996.PersonnelExample
\begin{example}\label{ex:semantics}
	Consider the instance $I$ of the
	relation \textit{Personnel} given in Table~\ref{tab:personnel},
	containing names of people, their position and their city (remember attribute
	names are not part of the model, they are just given here for ease of
	reading). Ignoring for now the~$t_i$'s, the
	relation has arity~$4$.
	Consider the following query: \emph{``What are the cities where at
		least two persons are working?''}. This
	can be expressed (without using aggregation) as:
	\(
		\qex\defeq\epsilon\left(\Pi_{\#4}\left(\textit{Personnel}\bowtie_{\#4=\#8\land\#1<\#5}
			\textit{Personnel}\right)\right).
	\)
	Note the need for duplicate elimination $\epsilon$ due to the
	multiset semantics. Then
  $\sem{I}{\qex}=\mset{(\textup{\revision{Nairobi}}),(\textup{Paris}),(\textup{\revision{Beijing}})}$.
\end{example}
% lax end

\begin{table}
	\centering
	\caption{The \textit{Personnel} relation
  (\revision{derived} from~\cite{dblp:journals/sigmod/senellart17})}\label{tab:personnel}
	\vspace{-.5em}
	\begin{tabular}{lllll}
		\toprule
		\underline{\textbf{id}} & \textbf{name} & \textbf{position} & \textbf{city}         \\
		\midrule
    1                       & \revision{Juma}          & Director          & \revision{Nairobi}      & $t_1$ \\
		2                       & Paul          & Janitor           & \revision{Nairobi}      & $t_2$ \\
    3                       & \revision{David}          & Analyst           & Paris         & $t_3$ \\
		4                       & Ellen         & Field agent       & \revision{Beijing}        & $t_4$ \\
    5                       & \revision{Aaheli}        & Double agent      & Paris         & $t_5$ \\
		6                       & Nancy         & HR                & Paris         & $t_6$ \\
    7                       & \revision{Jing}          & Analyst           & \revision{Beijing}        & $t_7$ \\
		\bottomrule
	\end{tabular}
\end{table}


\begin{toappendix}
	We also give SQL translations
	of the operators defined in this section, assuming appropriate
	translations of unnamed positional indices to attribute names within
	terms (when it makes a difference, we use
	PostgreSQL syntax variants):
	\begin{description}
		\item[relation]~
		\begin{minted}[escapeinside=||]{sql}
SELECT * FROM |$R$|
    \end{minted}
		\item[projection]~
		\begin{minted}[escapeinside=||]{sql}
SELECT |$t_1, t_2, \dots, t_n$| FROM (|$q$|)
    \end{minted}
		\item[selection]~
		\begin{minted}[escapeinside=||]{sql}
SELECT * FROM (|$q$|) WHERE |$\phi$|
    \end{minted}
		\item[cross product]~
		\begin{minted}[escapeinside=||]{sql}
SELECT * FROM (|$q_1$|), (|$q_2$|)
    \end{minted}
		\item[multiset sum]~
		\begin{minted}[escapeinside=||]{sql}
|$q_1$| UNION ALL |$q_2$|
    \end{minted}
		\item[duplicate elimination]~
		\begin{minted}[escapeinside=||]{sql}
SELECT DISTINCT * FROM (|$q$|)
    \end{minted}
		\item[multiset difference]~
		\begin{minted}[escapeinside=||]{sql}
SELECT * FROM (|$q_1$|) AS temp WHERE ROW(temp.*) NOT IN (|$q_2$|)
    \end{minted}
		\item[aggregation]~If $\#j$ stands for the name of the $j$th attribute
		returned by the query $q$:
		\begin{minted}[escapeinside=||]{sql}
SELECT |$\#i_1,\dots,\#i_m, f_1(t_1),\dots, f_n(t_n)$| FROM (|$q$|)
GROUP BY 1, |$\dots$|, m
    \end{minted}
		\item[join]~
		\begin{minted}[escapeinside=||]{sql}
SELECT * FROM (|$q_1$|) JOIN (|$q_2$|) ON |$\phi$|
    \end{minted}
		or simply
		\begin{minted}[escapeinside=||]{sql}
SELECT * FROM (|$q_1$|), (|$q_2$|) WHERE |$\phi$|
    \end{minted}
		\item[set union]~
		\begin{minted}[escapeinside=||]{sql}
|$q_1$| UNION |$q_2$|
    \end{minted}
	\end{description}

	In addition to relations being unlabeled and untyped, our query semantics slightly differs from SQL in the following way:
	\begin{compactitem}
		\item As mentioned, our semantics for multiset difference
		follows \mintinline{sql}{NOT IN}, not \mintinline{sql}{EXCEPT ALL}.
		\item We do not consider \mintinline{sql}{NULL} values in any way.
		\item We do not reproduce some quirks of SQL groupless aggregation: in SQL,
		\mintinline{sql}{SELECT COUNT(x) FROM R} returns 0 on an empty relation
		of schema \mintinline{sql}{R(x INT)} while \mintinline{sql}{SELECT SUM(x) FROM R}
		returns \mintinline{sql}{NULL}. In our semantics, $\gamma_{\#1}[\#1:
				f](R)$ returns the empty relation, however $f$ is defined, if evaluated on
		an instance where $R$ is empty.
		This is more consistent with
		aggregation under grouping.
	\end{compactitem}
\end{toappendix}
